\documentclass{ximera}

\input{../../../preamble.tex}


\definecolor{fvdnavy}{HTML}{071D43}
\definecolor{fvdcyan}{HTML}{20BFC6}
\definecolor{fvdcyanlight}{HTML}{EFFCFB}
\definecolor{fvdmagenta}{HTML}{F10D69}
\definecolor{fvdmagentalight}{HTML}{FFF1F7}
\definecolor{fvdtext}{HTML}{163044}





%=================================================
% Pedagogische omgevingen
% PDF: tcolorbox
% HTML: learning-box
%=================================================

\newenvironment{voorkennis}
{%
    \ifdefined\HCode
        \HCode{%
<section class="learning-box learning-box--prior">
<div class="learning-box__header">
<span class="learning-box__icon" aria-hidden="true"></span>
<span class="learning-box__title">Voorkennis</span>
</div>
<div class="learning-box__content">
        }%
        \begin{itemize}
    \else
       \begin{tcolorbox}[
    enhanced,
    breakable,
    colback=fvdcyanlight,
    colframe=fvdcyan,
    coltext=fvdtext,
    boxrule=0.5pt,
    leftrule=4pt,
    arc=4mm,
    outer arc=4mm,
    left=7mm,
    right=6mm,
    top=5mm,
    bottom=5mm,
    before skip=6mm,
    after skip=6mm,
    title=Voorkennis,
    coltitle=fvdcyan,
    fonttitle=\bfseries\Large,
    attach title to upper={\par\smallskip},
    boxed title style={
        empty,
        left=0mm,
        right=0mm,
        top=0mm,
        bottom=0mm
    },
    overlay={
        \draw[
            fvdcyan,
            line width=0.5pt
        ]
        ([xshift=42mm,yshift=-13mm]frame.north west)
        --
        ([xshift=-7mm,yshift=-13mm]frame.north east);
    }
]
        \begin{itemize}
    \fi
}
{%
    \end{itemize}
    \ifdefined\HCode
        \HCode{%
</div>
</section>
        }%
    \else
        \end{tcolorbox}
    \fi
}


\newenvironment{leerdoelen}
{%
    \ifdefined\HCode
        \HCode{%
<section class="learning-box learning-box--goals">
<div class="learning-box__header">
<span class="learning-box__icon" aria-hidden="true"></span>
<span class="learning-box__title">Leerdoelen</span>
</div>
<div class="learning-box__content">
        }%
        \begin{itemize}
    \else
\begin{tcolorbox}[
    enhanced,
    breakable,
    colback=fvdmagentalight,
    colframe=fvdmagenta,
    coltext=fvdtext,
    boxrule=0.5pt,
    leftrule=4pt,
    arc=4mm,
    outer arc=4mm,
    left=7mm,
    right=6mm,
    top=5mm,
    bottom=5mm,
    before skip=6mm,
    after skip=6mm,
    title=Leerdoelen,
    coltitle=fvdmagenta,
    fonttitle=\bfseries\Large,
    attach title to upper={\par\smallskip},
    boxed title style={
        empty,
        left=0mm,
        right=0mm,
        top=0mm,
        bottom=0mm
    },
    overlay={
        \draw[
            fvdmagenta,
            line width=0.5pt
        ]
        ([xshift=42mm,yshift=-13mm]frame.north west)
        --
        ([xshift=-7mm,yshift=-13mm]frame.north east);
    }
]
        \begin{itemize}
    \fi
}
{%
    \end{itemize}
    \ifdefined\HCode
        \HCode{%
</div>
</section>
        }%
    \else
        \end{tcolorbox}
    \fi
}


\addPrintStyle{../../..}

\title{Oefeningen — Primitieve functies van veeltermfuncties}
\author{Ingmar Herreman}

\begin{document}

% FVD_CHAPTER_DATA_START
% Automatisch gegenereerd uit index.tex.
% Niet handmatig aanpassen.
\fvdchapterdata{2}{Van oppervlakte naar integraal}{3}
% FVD_CHAPTER_DATA_END


% ============================================================
% OEFENING 1 — HERKENNEN EN CONTROLEREN
% ============================================================

\begin{oefening}[Is dit een primitieve?]{\oefB\ESS}

Gegeven is
\[
f(x)=6x^2-4x+3.
\]

\begin{question}
Controleer door af te leiden of
\[
F(x)=2x^3-2x^2+3x+7
\]
een primitieve functie van $f$ is.
\end{question}

\begin{question}
Controleer hetzelfde voor
\[
G(x)=2x^3-2x^2+3x-12.
\]
\end{question}

\begin{question}
Is
\[
H(x)=2x^3-2x^2+3x^2+1
\]
een primitieve van $f$? Motiveer.
\end{question}

\begin{question}
Leg uit waarom $F$ en $G$ slechts door een constante van elkaar verschillen.
\end{question}

\begin{oplossing}

We gebruiken de definitie: $F$ is een primitieve van $f$ als $F'(x)=f(x)$.

\begin{enumerate}
\item
\[
F'(x)=6x^2-4x+3=f(x).
\]
Dus $F$ is een primitieve van $f$.

\item
\[
G'(x)=6x^2-4x+3=f(x).
\]
Dus ook $G$ is een primitieve van $f$.

\item
\[
H'(x)=6x^2-4x+6x=6x^2+2x.
\]
Dit is niet gelijk aan $f(x)$, dus $H$ is geen primitieve van $f$.

\item
\[
F(x)-G(x)=19.
\]
Het verschil is constant. Dat is precies wat we verwachten voor twee primitieve functies van dezelfde functie.
\end{enumerate}

\end{oplossing}

\end{oefening}


% ============================================================
% OEFENING 2 — FAMILIE EN +C
% ============================================================

\begin{oefening}[Een hele familie primitieve functies]{\oefB\ESS}

Gegeven is
\[
f(x)=3x^2.
\]

\begin{question}
Toon aan dat elk van de volgende functies een primitieve is van $f$:
\[
F_1(x)=x^3,
\qquad
F_2(x)=x^3+5,
\qquad
F_3(x)=x^3-17.
\]
\end{question}

\begin{question}
Geef nog twee andere primitieve functies van $f$.
\end{question}

\begin{question}
Noteer de volledige familie primitieve functies met behulp van de onbepaalde integraal.
\end{question}

\begin{question}
Waarom mag de integratieconstante $C$ niet worden weggelaten wanneer naar de onbepaalde integraal wordt gevraagd?
\end{question}

\begin{oplossing}

\begin{enumerate}
\item
Voor alle drie geldt
\[
F_i'(x)=3x^2.
\]
Ze zijn dus allemaal primitieve functies van $f$.

\item
Bijvoorbeeld
\[
x^3+2
\qquad\text{en}\qquad
x^3-100.
\]

\item
\[
\boxed{\int 3x^2\,dx=x^3+C.}
\]

\item
Omdat er oneindig veel primitieve functies zijn. Ze verschillen onderling door een constante. Zonder $+C$ geef je slechts één lid van die familie.
\end{enumerate}

\end{oplossing}

\end{oefening}


% ============================================================
% OEFENING 3 — NOTATIE
% ============================================================

\begin{oefening}[Lees de onbepaalde integraal]{\oefB}

Beschouw
\[
\int (4x^3-2x+5)\,dx.
\]

\begin{question}
Wat is de integrand?
\end{question}

\begin{question}
Wat is de integratievariabele?
\end{question}

\begin{question}
Wat stelt het symbool $C$ voor in het antwoord?
\end{question}

\begin{question}
Leg in woorden het verschil uit tussen
\[
\int f(x)\,dx
\qquad\text{en}\qquad
\int_a^b f(x)\,dx.
\]
\end{question}

\begin{oplossing}

\begin{enumerate}
\item De integrand is $4x^3-2x+5$.
\item De integratievariabele is $x$.
\item $C$ is de integratieconstante en kan elke reële waarde aannemen.
\item De onbepaalde integraal beschrijft een familie primitieve functies. De bepaalde integraal is een getal en stelt een georiënteerde accumulatie over een interval voor.
\end{enumerate}

\end{oplossing}

\end{oefening}


% ============================================================
% OEFENING 4 — MACHTSREGEL
% ============================================================

\begin{oefening}[De machtsregel inoefenen]{\oefB\ESS}

Bepaal telkens de onbepaalde integraal.

\begin{question}
\[
\int x^2\,dx
\]
\end{question}

\begin{question}
\[
\int x^5\,dx
\]
\end{question}

\begin{question}
\[
\int x^8\,dx
\]
\end{question}

\begin{question}
\[
\int x\,dx
\]
\end{question}

\begin{question}
\[
\int 1\,dx
\]
\end{question}

\begin{question}
\[
\int x^{12}\,dx
\]
\end{question}

\begin{oplossing}

Met
\[
\int x^n\,dx=\frac{x^{n+1}}{n+1}+C
\qquad(n\neq -1)
\]
krijgen we:

\begin{enumerate}
\item $\displaystyle \frac{x^3}{3}+C$.
\item $\displaystyle \frac{x^6}{6}+C$.
\item $\displaystyle \frac{x^9}{9}+C$.
\item $\displaystyle \frac{x^2}{2}+C$.
\item $\displaystyle x+C$.
\item $\displaystyle \frac{x^{13}}{13}+C$.
\end{enumerate}

\end{oplossing}

\end{oefening}


% ============================================================
% OEFENING 5 — CONSTANTE FACTOREN
% ============================================================

\begin{oefening}[Constante factoren]{\oefB\ESS}

Bepaal telkens de volledige familie primitieve functies.

\begin{question}
\[
f(x)=4x^3.
\]
\end{question}

\begin{question}
\[
g(x)=-7x^5.
\]
\end{question}

\begin{question}
\[
h(x)=\frac{3}{2}x^4.
\]
\end{question}

\begin{question}
\[
p(x)=-\frac{5}{3}x^2.
\]
\end{question}

\begin{question}
\[
q(x)=9.
\]
\end{question}

\begin{oplossing}

\begin{enumerate}
\item $F(x)=x^4+C$.
\item $G(x)=-\dfrac{7}{6}x^6+C$.
\item $H(x)=\dfrac{3}{10}x^5+C$.
\item $P(x)=-\dfrac{5}{9}x^3+C$.
\item $Q(x)=9x+C$.
\end{enumerate}

\end{oplossing}

\end{oefening}


% ============================================================
% OEFENING 6 — VEELTERMFUNCTIES, OUDE REEKS GEÏNTEGREERD
% ============================================================

\begin{oefening}[Veeltermfuncties term voor term integreren]{\oefB\ESS}

Bepaal telkens de volledige familie primitieve functies.

\begin{question}
\[
f(x)=-\frac{4}{3}x^3+\frac{2}{3}x.
\]
\end{question}

\begin{question}
\[
g(x)=2x-4.
\]
\end{question}

\begin{question}
\[
h(x)=\frac12x^3+6x^2-4x+8.
\]
\end{question}

\begin{question}
\[
p(x)=-\frac23x^2+\frac34x.
\]
\end{question}

\begin{question}
\[
q(x)=8x^3-12x^2+6x-1.
\]
\end{question}

\begin{oplossing}

\begin{enumerate}
\item
\[
F(x)=-\frac{x^4}{3}+\frac{x^2}{3}+C.
\]

\item
\[
G(x)=x^2-4x+C.
\]

\item
\[
H(x)=\frac{x^4}{8}+2x^3-2x^2+8x+C.
\]

\item
\[
P(x)=-\frac{2x^3}{9}+\frac{3x^2}{8}+C.
\]

\item
\[
Q(x)=2x^4-4x^3+3x^2-x+C.
\]
\end{enumerate}

\end{oplossing}

\end{oefening}


% ============================================================
% OEFENING 7 — MEER VEELTERMEN
% ============================================================

\begin{oefening}[Veeltermfuncties — extra inoefening]{\oefB}

Bepaal telkens de onbepaalde integraal.

\begin{question}
\[
\int (6x^5-3x^2+8x-4)\,dx
\]
\end{question}

\begin{question}
\[
\int (5x^4+2x^3-7x+9)\,dx
\]
\end{question}

\begin{question}
\[
\int (-3x^6+4x^2-5)\,dx
\]
\end{question}

\begin{question}
\[
\int \left(\frac34x^3-\frac25x^2+\frac12x-6\right)\,dx
\]
\end{question}

\begin{oplossing}

\begin{enumerate}
\item
\[
x^6-x^3+4x^2-4x+C.
\]
\item
\[
x^5+\frac12x^4-\frac72x^2+9x+C.
\]
\item
\[
-\frac37x^7+\frac43x^3-5x+C.
\]
\item
\[
\frac{3}{16}x^4-\frac{2}{15}x^3+\frac14x^2-6x+C.
\]
\end{enumerate}

\end{oplossing}

\end{oefening}


% ============================================================
% OEFENING 8 — EERST ALGEBRA
% ============================================================

\begin{oefening}[Eerst vereenvoudigen, dan integreren]{\oefB\ESS}

Vereenvoudig eerst de integrand. Bepaal daarna de onbepaalde integraal.

\begin{question}
\[
\int \left(2x(x^2-3)+x^3+4\right)\,dx
\]
\end{question}

\begin{question}
\[
\int \left((x+2)(x^2-3x+1)\right)\,dx
\]
\end{question}

\begin{question}
\[
\int \left(3x^2-2(x^2-4x+1)+5\right)\,dx
\]
\end{question}

\begin{question}
\[
\int \left(x(x-1)(x+1)+2x\right)\,dx
\]
\end{question}

\begin{oplossing}

\begin{enumerate}
\item
\[
2x(x^2-3)+x^3+4=3x^3-6x+4,
\]
dus
\[
\boxed{\frac34x^4-3x^2+4x+C}.
\]

\item
\[
(x+2)(x^2-3x+1)=x^3-x^2-5x+2,
\]
dus
\[
\boxed{\frac14x^4-\frac13x^3-\frac52x^2+2x+C}.
\]

\item
\[
3x^2-2(x^2-4x+1)+5=x^2+8x+3,
\]
dus
\[
\boxed{\frac13x^3+4x^2+3x+C}.
\]

\item
\[
x(x-1)(x+1)+2x=x(x^2-1)+2x=x^3+x,
\]
dus
\[
\boxed{\frac14x^4+\frac12x^2+C}.
\]
\end{enumerate}

\end{oplossing}

\end{oefening}


% ============================================================
% OEFENING 9 — CONTROLEREN DOOR AFLEIDEN
% ============================================================

\begin{oefening}[Controleer je antwoord]{\oefB\ESS}

Een leerling heeft telkens een onbepaalde integraal berekend. Controleer het antwoord door af te leiden. Verbeter wanneer nodig.

\begin{question}
\[
\int 4x^3\,dx=4x^4+C.
\]
\end{question}

\begin{question}
\[
\int (6x^2-4x+5)\,dx=2x^3-2x^2+5x+C.
\]
\end{question}

\begin{question}
\[
\int (3x^4+2)\,dx=\frac35x^5+2x+C.
\]
\end{question}

\begin{question}
\[
\int (-8x+3)\,dx=-4x^2+3x+C.
\]
\end{question}

\begin{oplossing}

\begin{enumerate}
\item Fout. $(4x^4+C)'=16x^3$. Correct is
\[
\boxed{\int 4x^3\,dx=x^4+C}.
\]

\item Correct, want
\[
(2x^3-2x^2+5x+C)'=6x^2-4x+5.
\]

\item Correct, want
\[
\left(\frac35x^5+2x+C\right)'=3x^4+2.
\]

\item Correct, want
\[
(-4x^2+3x+C)'=-8x+3.
\]
\end{enumerate}

\end{oplossing}

\end{oefening}


% ============================================================
% OEFENING 10 — +C EN REDENEREN
% ============================================================

\begin{oefening}[Waarom $+C$?]{\oefU\ESS}

Een leerling schrijft
\[
\int (6x-2)\,dx=3x^2-2x.
\]

\begin{question}
Is $3x^2-2x$ een primitieve van $6x-2$? Leg uit.
\end{question}

\begin{question}
Waarom is het antwoord toch onvolledig als gevraagd wordt naar de onbepaalde integraal?
\end{question}

\begin{question}
Schrijf de correcte onbepaalde integraal.
\end{question}

\begin{question}
Zijn de functies
\[
3x^2-2x+4
\qquad\text{en}\qquad
3x^2-2x-\pi
\]
ook primitieve functies van $6x-2$?
\end{question}

\begin{oplossing}

\begin{enumerate}
\item Ja, want $(3x^2-2x)'=6x-2$.
\item De onbepaalde integraal vraagt naar de volledige familie primitieve functies.
\item
\[
\boxed{\int(6x-2)\,dx=3x^2-2x+C.}
\]
\item Ja. Constante termen verdwijnen bij het afleiden.
\end{enumerate}

\end{oplossing}

\end{oefening}


% ============================================================
% OEFENING 11 — SPECIFIEKE PRIMITIEVE: OUDE OPGAVE + EXTRA
% ============================================================

\begin{oefening}[Een primitieve bepalen met een voorwaarde]{\oefB\ESS}

Bepaal telkens de primitieve functie $F$ die voldoet aan de gegeven voorwaarde.

\begin{question}
\[
f(x)=8x^2-x^2+4x-3,
\qquad
F(1)=2.
\]
\end{question}

\begin{question}
\[
f(x)=6x^2-4x+3,
\qquad
F(1)=5.
\]
\end{question}

\begin{question}
\[
f(x)=4x^3-2x,
\qquad
F(0)=-3.
\]
\end{question}

\begin{question}
\[
f(x)=3x^2+2x-1,
\qquad
F(-1)=4.
\]
\end{question}

\begin{oplossing}

\begin{enumerate}
\item Eerst
\[
f(x)=7x^2+4x-3.
\]
Dus
\[
F(x)=\frac73x^3+2x^2-3x+C.
\]
Uit $F(1)=2$ volgt
\[
\frac73+2-3+C=2
\]
en dus
\[
C=\frac23.
\]
Daarom
\[
\boxed{F(x)=\frac73x^3+2x^2-3x+\frac23}.
\]

\item
\[
F(x)=2x^3-2x^2+3x+C.
\]
Uit $F(1)=5$:
\[
2-2+3+C=5,
\]
dus $C=2$. Daarom
\[
\boxed{F(x)=2x^3-2x^2+3x+2}.
\]

\item
\[
F(x)=x^4-x^2+C.
\]
Uit $F(0)=-3$ volgt $C=-3$. Dus
\[
\boxed{F(x)=x^4-x^2-3}.
\]

\item
\[
F(x)=x^3+x^2-x+C.
\]
Uit $F(-1)=4$ volgt
\[
-1+1+1+C=4,
\]
dus $C=3$. Daarom
\[
\boxed{F(x)=x^3+x^2-x+3}.
\]
\end{enumerate}

\end{oplossing}

\end{oefening}


% ============================================================
% OEFENING 12 — VOORWAARDEN EN PARAMETER
% ============================================================

\begin{oefening}[Welke constante hoort erbij?]{\oefU}

Gegeven is
\[
f(x)=12x^2-6x+4.
\]

\begin{question}
Bepaal de volledige familie primitieve functies van $f$.
\end{question}

\begin{question}
Bepaal de primitieve $F$ waarvoor $F(0)=7$.
\end{question}

\begin{question}
Bepaal de primitieve $G$ waarvoor $G(2)=0$.
\end{question}

\begin{question}
Bepaal het constante verschil $F(x)-G(x)$.
\end{question}

\begin{oplossing}

\begin{enumerate}
\item
\[
4x^3-3x^2+4x+C.
\]

\item Uit $F(0)=7$ volgt $C=7$, dus
\[
F(x)=4x^3-3x^2+4x+7.
\]

\item
\[
G(2)=32-12+8+C=28+C=0,
\]
dus $C=-28$ en
\[
G(x)=4x^3-3x^2+4x-28.
\]

\item
\[
F(x)-G(x)=35.
\]
\end{enumerate}

\end{oplossing}

\end{oefening}


% ============================================================
% OEFENING 13 — NEGATIEVE EN RATIONALE EXPONENTEN UIT OUDE REEKS
% ============================================================

\begin{oefening}[Eerst herschrijven als een macht]{\oefU}

De volgende functies zijn geen veeltermfuncties. Schrijf ze eerst met machten en pas daarna de machtsregel toe. Gebruik de regel alleen wanneer de exponent niet gelijk is aan $-1$.

\begin{question}
\[
f(x)=\frac{1}{x^2}-1+2x^2.
\]
\end{question}

\begin{question}
\[
g(x)=\frac{1}{\sqrt{x}}+2.
\]
\end{question}

\begin{question}
\[
h(x)=\sqrt[3]{x^5}.
\]
\end{question}

\begin{question}
\[
p(x)=\frac{1}{\sqrt{x}}.
\]
\end{question}

\begin{question}
\[
q(x)=\frac{2}{x^3}+4.
\]
\end{question}

\begin{question}
\[
r(x)=\frac{1}{x^3}+1+x^2.
\]
\end{question}

\begin{oplossing}

We gebruiken
\[
\frac1{x^2}=x^{-2},
\qquad
\frac1{\sqrt{x}}=x^{-1/2},
\qquad
\sqrt[3]{x^5}=x^{5/3}.
\]

\begin{enumerate}
\item
\[
F(x)=-\frac1x-x+\frac23x^3+C.
\]

\item
\[
G(x)=2\sqrt{x}+2x+C.
\]

\item
\[
H(x)=\frac38x^{8/3}+C
=\frac38\sqrt[3]{x^8}+C.
\]

\item
\[
P(x)=2\sqrt{x}+C.
\]

\item
\[
Q(x)=4x-\frac1{x^2}+C.
\]

\item
\[
R(x)=-\frac1{2x^2}+x+\frac{x^3}{3}+C.
\]
\end{enumerate}

\end{oplossing}

\end{oefening}


% ============================================================
% OEFENING 14 — F EN F' RELATIE
% ============================================================

\begin{oefening}[Wat vertelt $f$ over haar primitieve?]{\oefU\ESS}

Een functie $F$ is een primitieve van
\[
f(x)=x^2-4.
\]

\begin{question}
Bepaal de nulpunten van $f$.
\end{question}

\begin{question}
Onderzoek het teken van $f$.
\end{question}

\begin{question}
Op welke intervallen is elke primitieve $F$ stijgend?
\end{question}

\begin{question}
Op welk interval is elke primitieve $F$ dalend?
\end{question}

\begin{question}
Wat kun je zeggen over de raaklijnen aan de grafiek van $F$ bij $x=-2$ en $x=2$?
\end{question}

\begin{question}
Moet je voor deze vragen eerst een formule voor $F$ bepalen? Leg uit.
\end{question}

\begin{oplossing}

Omdat $F'(x)=f(x)$, bepaalt het teken van $f$ het verloop van $F$.

\begin{enumerate}
\item
\[
x^2-4=0
\quad\Longleftrightarrow\quad
x=-2\text{ of }x=2.
\]

\item
\[
f(x)>0 \text{ voor }x<-2\text{ en }x>2,
\]
terwijl
\[
f(x)<0 \text{ voor }-2<x<2.
\]

\item $F$ is stijgend op $]-\infty,-2[$ en $]2,+\infty[$.
\item $F$ is dalend op $]-2,2[$.
\item Omdat $F'(-2)=F'(2)=0$, zijn de raaklijnen daar horizontaal.
\item Nee. De relatie $F'=f$ volstaat om het verloop te onderzoeken.
\end{enumerate}

\end{oplossing}

\end{oefening}


% ============================================================
% OEFENING 15 — OMGEKEERD REDENEREN
% ============================================================

\begin{oefening}[Van primitieve naar integrand]{\oefB}

Gegeven is een functie $F$. Bepaal telkens de functie $f$ waarvoor $F$ een primitieve is.

\begin{question}
\[
F(x)=2x^4-3x^2+7x-5.
\]
\end{question}

\begin{question}
\[
F(x)=-\frac12x^6+4x^3-2.
\]
\end{question}

\begin{question}
\[
F(x)=\frac35x^5-\frac23x^3+x+9.
\]
\end{question}

\begin{oplossing}

We leiden $F$ af.

\begin{enumerate}
\item
\[
f(x)=8x^3-6x+7.
\]
\item
\[
f(x)=-3x^5+12x^2.
\]
\item
\[
f(x)=3x^4-2x^2+1.
\]
\end{enumerate}

\end{oplossing}

\end{oefening}


% ============================================================
% OEFENING 16 — FOUTENANALYSE
% ============================================================

\begin{oefening}[Zoek de fout]{\oefU\ESS}

Onderzoek elke redenering. Geef aan wat fout loopt en verbeter het antwoord.

\begin{question}
Een leerling schrijft
\[
\int x^4\,dx=4x^3+C.
\]
\end{question}

\begin{question}
Een leerling schrijft
\[
\int 5x^2\,dx=\frac53x^3.
\]
\end{question}

\begin{question}
Een leerling schrijft
\[
\int (2x^3-4x)\,dx
=\frac12x^4-2x^2+C.
\]
\end{question}

\begin{question}
Een leerling beweert: ``Als $F$ en $G$ beide primitieve functies van $f$ zijn, dan zijn $F$ en $G$ noodzakelijk gelijk.''
\end{question}

\begin{oplossing}

\begin{enumerate}
\item De machtsregel voor afgeleiden werd ten onrechte gebruikt. Correct:
\[
\boxed{\int x^4\,dx=\frac15x^5+C}.
\]

\item De primitieve zelf is correct, maar $+C$ ontbreekt. Correct:
\[
\boxed{\int 5x^2\,dx=\frac53x^3+C}.
\]

\item Correct, want
\[
\left(\frac12x^4-2x^2+C\right)'=2x^3-4x.
\]

\item Fout. Twee primitieve functies van dezelfde functie kunnen door een constante verschillen:
\[
F(x)=G(x)+C.
\]
\end{enumerate}

\end{oplossing}

\end{oefening}


% ============================================================
% OEFENING 17 — GEMENGDE ZELFCHECK
% ============================================================

\begin{oefening}[Zelfcheck: primitieve functies]{\oefU\ESS}

Gegeven is
\[
f(x)=4x^3-6x^2+2x.
\]

\begin{question}
Bepaal de volledige familie primitieve functies van $f$.
\end{question}

\begin{question}
Bepaal de primitieve $F$ waarvoor $F(0)=-2$.
\end{question}

\begin{question}
Controleer je antwoord door $F$ af te leiden.
\end{question}

\begin{question}
Bepaal de punten waar de grafiek van $F$ een horizontale raaklijn kan hebben. Je hoeft de functiewaarden van die punten niet te bepalen.
\end{question}

\begin{question}
Leg uit welk verband je bij de vorige vraag gebruikt tussen $f$ en $F$.
\end{question}

\begin{oplossing}

\begin{enumerate}
\item
\[
\int (4x^3-6x^2+2x)\,dx
=x^4-2x^3+x^2+C.
\]

\item Uit $F(0)=-2$ volgt $C=-2$. Dus
\[
\boxed{F(x)=x^4-2x^3+x^2-2}.
\]

\item
\[
F'(x)=4x^3-6x^2+2x=f(x).
\]

\item Horizontale raaklijnen voldoen aan $F'(x)=0$. Omdat $F'=f$, moeten we dus
\[
4x^3-6x^2+2x=0.
\]
We ontbinden:
\[
2x(2x^2-3x+1)=2x(2x-1)(x-1)=0.
\]
Dus
\[
x=0,\qquad x=\frac12,\qquad x=1.
\]
Bij deze drie abscissen heeft de grafiek van $F$ een horizontale raaklijn.

\item We gebruiken
\[
F'(x)=f(x).
\]
Een horizontale raaklijn aan $F$ ontstaat waar $F'(x)=0$, dus waar $f(x)=0$.
\end{enumerate}

\end{oplossing}

\end{oefening}


% ============================================================
% OEFENING 18 — GEVORDERD: REDENEREN ZONDER ALLES UIT TE REKENEN
% ============================================================

\begin{oefening}[Primitieven vergelijken]{\oefG}

Van twee functies $F$ en $G$ weten we dat
\[
F'(x)=G'(x)=2x^3-6x.
\]
Verder geldt
\[
F(1)=4
\qquad\text{en}\qquad
G(1)=-3.
\]

\begin{question}
Wat kun je besluiten over $F(x)-G(x)$ zonder de functies volledig te bepalen?
\end{question}

\begin{question}
Bepaal de waarde van $F(x)-G(x)$.
\end{question}

\begin{question}
Stel dat $H$ nog een andere primitieve is van $2x^3-6x$. Welke vorm heeft $H$ ten opzichte van $F$?
\end{question}

\begin{question}
Bepaal nu ter controle expliciet de algemene primitieve van $2x^3-6x$.
\end{question}

\begin{oplossing}

\begin{enumerate}
\item Omdat $F'$ en $G'$ gelijk zijn, is
\[
(F-G)'=0.
\]
Dus $F-G$ is constant.

\item Bij $x=1$:
\[
F(1)-G(1)=4-(-3)=7.
\]
Dus
\[
\boxed{F(x)-G(x)=7.}
\]

\item Er bestaat een constante $k$ zodat
\[
H(x)=F(x)+k.
\]

\item
\[
\int (2x^3-6x)\,dx
=\frac12x^4-3x^2+C.
\]
\end{enumerate}

\end{oplossing}

\end{oefening}

\end{document}

